\documentclass[10pt]{beamer}
% \documentclass[handout, 10pt]{beamer}

\usepackage{homework}

\usetheme{Madrid}
\setbeamersize{text margin left=30pt,text margin right=20pt}
\graphicspath{{../../figures/}}

\usepackage[backend=biber,style=authoryear-comp]{biblatex}
\addbibresource{references.bib}
% Beamer needs this to show footnotes properly
\setbeamertemplate{footnote}{%
  \parindent 0em\noindent
  \raggedright
  \insertfootnotemark\insertfootnotetext\par
}
% Change item symbols for bibliography
\setbeamertemplate{bibliography item}[triangle]
\setbeamertemplate{itemize item}[triangle]

% Title Page
\title[Simulating Sandpiles]{Simulating Sandpiles}
\subtitle{Computational Physics (physics760)}
\author{Konrad Beck, Leon Galbas}
\date{09.03.2026}
\institute[Universität Bonn]{Rheinische Friedrich-Wilhelms-Universität Bonn}


\begin{document}
\frame{\titlepage}

% Table of Contents
\begin{frame}{Outline}
    \tableofcontents
\end{frame}
% Current section
\AtBeginSection[ ]
{
\begin{frame}{Outline}
    \tableofcontents[currentsection]
\end{frame}
}
\AtBeginSubsection[]
{
  \begin{frame}{Outline}
    \tableofcontents[currentsection,currentsubsection]
  \end{frame}
}

% Start of actual presentation

\section{Motivation}
\begin{frame}{Introduction}

    \begin{figure}
        \centering

        \begin{subfigure}{0.48\textwidth}
            \centering
            \includegraphics[width=\linewidth,height=0.55\textheight,keepaspectratio]{ising_example_clustersize.png}
            \caption{Log-log plot of the cluster-size distribution $n_I(s,V)$ versus $s$ for an Ising model in the Ising sector \cite{fang2021percolation}.}
        \end{subfigure}
        \hfill
        \begin{subfigure}{0.48\textwidth}
            \centering
            \includegraphics[width=\linewidth,height=0.55\textheight,keepaspectratio]{earthquake_power_law}
            \caption{Earthquake magnitude distribution in Southern California showing a power-law behavior over 6 decades \cite{bak2002unified}.}
        \end{subfigure}

        \caption{Two different systems in a critical state showing behavior governed by power laws.}
    \end{figure}

    Question: What is the difference between these two?

\end{frame}



\begin{frame}{Self-organized criticality}
    Answer: The Ising model needs a fine-tuned parameter (e.g. temperature) to reach criticality, while earthquakes exhibit this behavior in their own. \\
        $\implies$ Self-Organized Criticality (SOC) \pause

    \begin{itemize}
        \item SOC naturally explains the widespread appearance of fractal structures and inverse-frequency noise in nature \pause
        \item Scale invariance: events of all sizes happen, no single "characteristic" scale
            $\implies$ Probability distributions of events follow power laws (Pareto distribution)
    \end{itemize}

\end{frame}



%\begin{frame}{Power laws}

    % TODO: Image of one or two power laws, e.g. solar flares (see https://arxiv.org/pdf/2212.02678)

%\end{frame}



\begin{frame}{Why sandpiles?}

    \begin{itemize}
        \item The Bak-Tang-Wiesenfeld (BTW) sandpile model is a very simple model to study SOC (\cite{bak1987self})
        \item Analogy: Slowly dropping sand onto a table causes regular avalanches of various sizes
    \end{itemize}

    \begin{figure}
        \centering
        \includegraphics[width=\linewidth,height=0.55\textheight,keepaspectratio]{sandpile_example}
        \caption{Example state of a two dimensional BTW sandpile model.}
    \end{figure}

\end{frame}



\begin{frame}{Goals of our project}

    \begin{itemize}
        \item Simulate a BTW sandpile cellular automaton in $d=2,3,4$ dimensions
        \item Show that the system exhibits self-organized criticality
        \item Confirm the emergence of power laws and compute the scaling exponents
        \item Compare the results for different hyperparameters and dimensions
    \end{itemize}

\end{frame}



\section{Sandpiles}
\begin{frame}{The BTW sandpile}
    \textbf{Definition:}

    \begin{itemize}
        \item A $d$-dimensional hypercubic grid $R\subset\bZ^d$ of linear size $N$.
        \item A function $z(r):\; \bZ^d \to \bZ$. (The local \enquote{slope} or \enquote{energy}.)
        \item A threshold $z_c$, s.t. for $z(r)>z_c$ the system is unstable.
        \item A relaxation function $z_\tau(r)\mapsto z_{\tau+1}(r)$.
    \end{itemize}
\end{frame}


\begin{frame}{The relaxation function}
    For all $r\in R$ with $z_\tau(r) > z_c$, i.e. at every unstable lattice site:
    \begin{align*}
        z_\tau(r) &\to z_\tau(r) - 2d, \\
        z_\tau(r\pm e_i) &\to z_\tau(r\pm e_i) + 1 \qq{for} i=1,\ldots,d.
    \end{align*}
    Imagine sandpile \enquote{toppling} onto its neighbors.
\end{frame}


\begin{frame}{An avalanche}
    \centering
    %\animategraphics[loop,width=\textwidth]{1}{../../figures/heatmaps/relaxation_example/frame-}{0}{7}
    % TODO: Insert animation of sandpile that starts of in a very unstable state and topples toward stable state (without further perturbation)
\end{frame}


\begin{frame}{Boundary conditions}

    \textbf{Open:}
    \begin{align*}
        z_\tau(r) &\to z_\tau(r) - 2d + (\text{number of }i\text{ with }r_i=N), \\
        z_\tau(r+e_i) &\to z_\tau(r+e_i) + 1 \qq{if} r_i\neq N, \\
        z_\tau(r-e_i) &\to z_\tau(r-e_i) + 1 \qq{for} i=1,\ldots,d, \\
        \text{and} z_\tau(r) &= 0 \qq{if} r_i=1 \text{ for any i}.
    \end{align*} \pause

    \textbf{Closed:}
    \begin{equation*}
        z_\tau(r) = 0 \qq{if} r_i=1\text{ or } r_i=N \text{ for any i}.
    \end{equation*}

\end{frame}


\begin{frame}{Driving the system}
    \textbf{Conservative perturbation:} Keeps the total energy constant.
    \begin{align*}
        z(r_0) &\to z(r_0) + d, \\
        z(r_0-e_i) &\to z(r_0-e_i) - 1 \qq{for} i=1,\ldots,d
    \end{align*}
    Local rearrangement of the sandpile. \pause

    \textbf{Nonconservative perturbation:} Locally increases the energy.
    \begin{equation*}
        z(r_0) \to z(r_0) + 1
    \end{equation*}
    Randomly drop a grain of sand.
\end{frame}


\begin{frame}{Driving the system}

    % TODO: Timelapse animation of empty system being driven until some avalanches happen.
    
    \centering
    %\animategraphics[width=\textwidth]{15}{../../figures/heatmaps/burn_in/frame-}{0}{499}
\end{frame}




\begin{frame}{Reaching criticality}

    \begin{itemize}
        \item Perturbation: Introduces energy to the system
        \item Relaxation: Dissipates energy from the system at the boundaries
        \item In the stationary state: \begin{equation*}
            \expval{z} \to \overline{\expval{z}} = \lim_{T\to\infty}\frac{1}{T}\int_{0}^{T}\expval{z}(\tau)\dd{\tau}.
        \end{equation*}
    \end{itemize}

    In the critical state, the total (spatial average) energy fluctuates around a temporal average energy.

\end{frame}


\begin{frame}{The critical state}

    % TODO: Animation of a system already in stationary state, exhibiting avalanches of various sizes
    \centering
    %\animategraphics[width=\textwidth]{15}{../../figures/heatmaps/stationary/frame-}{0}{499}
\end{frame}


\begin{frame}{What do we measure?}

    \textbf{Avalanche variables:}
    \begin{itemize}
        \item Avalanche lifetime $t$.
        \item Total size $s$. \\ (Cumulative sum of all toppling events.)
        \item Spatial linear size $l$. \\ (Maximum linear extend in any dimension.)
    \end{itemize}

\end{frame}


\begin{frame}{What do we expect?}

    \textbf{Power laws:}
    \begin{align*}
        P(T=t) &\propto t^{1-\alpha}, \\
        P(S=s) &\propto s^{1-\tau}, \\
        P(L=l) &\propto l^{1-\lambda}.
    \end{align*}

\end{frame}



\section{Implementation}
\begin{frame}{The transient phase}

    % Question (spoken): How do we know whether a system has reached a stationary state
    \pause

    \begin{figure}
        \centering

        \begin{subfigure}{0.48\textwidth}
            \centering
            \includegraphics[width=\linewidth,height=0.55\textheight,keepaspectratio]{N40d2_open_conservative_zero_burn_in}
            \caption{Conservative perturbation}
        \end{subfigure}
        \hfill
        \begin{subfigure}{0.48\textwidth}
            \centering
            \includegraphics[width=\linewidth,height=0.55\textheight,keepaspectratio]{N40d2_open_nonconservative_zero_burn_in}
            \caption{Nonconservative perturbation}
        \end{subfigure}

        \caption{The spatial average $\expval{z}(\tau)$ over time for two sandpiles of dimension $N=40,\,d=2$ initialized to all zeroes with open boundary conditions.}
    \end{figure}
\end{frame}


\begin{frame}{How to detect stationarity?}
    \pause
    \begin{enumerate}
        \item Initialize the sandpile.
        \item Define a threshold $\eps>0$.
        \item Track $\expval{z}$ every $\num{1000}$ steps while continually perturbing and relaxing the sandpile.
        \item Do a linear fit through the last $\num{50}$ values.
        \item Stop, if the slope is smaller than $\eps$.
    \end{enumerate}
\end{frame}


\begin{frame}{The linear scaling region}

    \begin{figure}
        \centering
        \includegraphics[width=\linewidth,height=0.55\textheight,keepaspectratio]{N20d4_open_conservative_alpha_no-boundaries}
        \caption{Measured probability distribution of $t$ for a sandpile of dimension $N=20,\,d=4$ with open boundaries and conservative perturbation.}
    \end{figure}

\end{frame}


\begin{frame}{Determining the scaling region}
    \begin{enumerate}
        \item Apply linear regression on rolling window of fixed size.
        \item Determine the slope $m_i$ and uncertainty $\Delta m_i$ for each window.
        \item Calculate differences\begin{equation*}
            d_i = \abs{m_i-m_{i+1}}, \qquad
            \Delta d_i = \sqrt{(\Delta m_i)^2+(\Delta m_{i+1})^2}.
        \end{equation*}
        \item Classify as stable, if $d_i < k\cdot\Delta d_i$, and if both fits satisfy $r_i^2 \geq r_{\mathrm{thresh}}$.
        \item Set upper and lower bounds of scaling region.
    \end{enumerate}
\end{frame}


\begin{frame}{Determining the scaling region}

    \begin{figure}
        \centering
        \includegraphics[width=\linewidth,height=0.55\textheight,keepaspectratio]{N20d4_open_conservative_alpha}
        \caption{Measured probability distribution of $t$ for a sandpile of dimension $N=20,\,d=4$ with open boundaries and conservative perturbation. The scaling region (fitting window) was chosen by the aforementioned heuristic.}
    \end{figure}

\end{frame}

\begin{frame}{Error estimation of scaling exponents}
    \pause 
    \begin{enumerate}
        \item compute upper and lower bounds of scaling region
        \item vary the boundaries by $\pm 3$
        \item compute the slope of every valid pair of boundaries
        \item compute standard deviation  $\sigma_\f{sys}$ of all slopes
        \item combine with linear fit error $\sigma_\f{stat}$
    \end{enumerate}
    \pause 

    \begin{equation*}
    \sigma_\f{tot}=\sqrt{\sigma_\f{stat}^{2}+\sigma_\f{sys}^{2}}
    \end{equation*}
\end{frame}



\section{Results}
\begin{frame}{What we measured}
    \begin{table}[ht!]
  \centering
  \begin{tabular}{lccccc}
    \hline
    config & $d$ & $N$ & boundary condition & perturbation & $n$ \\
    \hline
    A & 2 & 40 & open   & nonconservative & $10^7$ \\
    B & 2 & 40 & open   & conservative    & $10^7$ \\
    C & 2 & 40 & closed & nonconservative & $10^7$ \\
    D & 3 & 20 & open   & nonconservative & $10^7$ \\
    E & 3 & 20 & open   & conservative    & $10^7$ \\
    F & 3 & 20 & closed & nonconservative & $10^7$ \\
    G & 4 & 20 & open   & nonconservative & $10^6$ \\
    H & 4 & 20 & open   & conservative    & $10^6$ \\
    I & 4 & 20 & closed & nonconservative & $10^6$ \\
    \hline
  \end{tabular}
  \caption{Simulated model configurations and measurement counts used in this work.}
  \label{tab:model_configurations}
\end{table}
\end{frame}

\begin{frame}{Did we observe power laws ?}
    \pause
    \begin{figure}
        \centering
        \includegraphics[width=\linewidth,height=0.55\textheight,keepaspectratio]{../../figures/loglog_plots/N40d2_open_conservative_lambda}
        \caption{Measured probability distribution of $l$ for a sandpile of dimension $N=40,\,d=2$ with open boundaries and conservative perturbation.}
    \end{figure}
\end{frame}

\begin{frame}{Did we observe power laws ?}

    \begin{figure}
        \centering
        \includegraphics[width=\linewidth,height=0.55\textheight,keepaspectratio]{../../figures/loglog_plots/N40d2_open_conservative_tau}
        \caption{Measured probability distribution of $s$ for a sandpile of dimension $N=40,\,d=2$ with open boundaries and conservative perturbation.}
    \end{figure}
\end{frame}


\begin{frame}{Did we observe power laws ?}
 
    \begin{figure}
        \centering
        \includegraphics[width=\linewidth,height=0.55\textheight,keepaspectratio]{../../figures/loglog_plots/N40d2_open_conservative_alpha}
        \caption{Measured probability distribution of $t$ for a sandpile of dimension $N=40,\,d=2$ with open boundaries and conservative perturbation.}
    \end{figure}
\end{frame}



\begin{frame}{Resulting exponentes}
    \begin{table}[!ht]
    \centering
    \begin{tabular}{cccc}
\toprule
model & $\tau$ & $\alpha$ & $\lambda$ \\
\midrule
A & $ (2.0182\pm0.0077) $ & $ (2.058\pm0.035) $ & $ (1.61\pm0.15) $ \\
D & $ (2.287\pm0.023) $ & $ (2.466\pm0.087) $ & $ (2.27\pm0.27) $ \\
G & $ (2.397\pm0.032) $ & $ (2.62\pm0.13) $ & $ (2.83\pm0.51) $ \\
\bottomrule
\end{tabular}

    \caption{Scaling exponents of the probability density functions for the open-boundary, nonconservative sandpile configurations (A, D, G). The exponents $\alpha$, $\tau$, and $\lambda$ are obtained from weighted linear least-squares fits on double-logarithmic PDF plots in the selected scaling window; The errors were determined by the aforementioned procedure.}
    \label{tab:scaling_exp_open_nonconservative}
    \end{table}
\end{frame}

\begin{frame}{Comparing the exponentes to literature}
    \begin{table}[!ht]
    \centering
    \begin{tabular}{cccc}
\toprule
model & $\tau$ & $\alpha$ & $\lambda$ \\
\midrule
A & $ (2.0182\pm0.0077) $ & $ (2.058\pm0.035) $ & $ (1.61\pm0.15) $ \\
D & $ (2.287\pm0.023) $ & $ (2.466\pm0.087) $ & $ (2.27\pm0.27) $ \\
G & $ (2.397\pm0.032) $ & $ (2.62\pm0.13) $ & $ (2.83\pm0.51) $ \\
\bottomrule
\end{tabular}

    \end{table}
    \begin{figure}
        \centering
        \includegraphics[width=\linewidth,height=0.55\textheight,keepaspectratio]{values_from_paper}
        \caption*{\cite{christensen1991dynamical}}
    \end{figure}
\end{frame} 




% Sources
\section*{Sources}
\begin{frame}{Sources: Software}
    The simulations and analyses described in this talk were implemented in Python, using the PyTorch and NumPy packages for efficient handling of multi dimensional data. Visualization was done in Matplotlib. The source code used to generate all numerical results is publicly available on GitHub at:

    \href{https://github.com/leon-galbas/ws2025-physics760-sandpile-simulation}{\textcolor{blue}{https://github.com/leon-galbas/ws2025-physics760-sandpile-simulation}}.
\end{frame}
\begin{frame}{Sources: Literature}
    \begin{tiny}
        \printbibliography
    \end{tiny}
\end{frame}

\end{document}
